% year: תשע"ח
% type: מבחן
% semester: א
% moed: ב
% number: 1ב
% points: 10
% categories: func-limits, lhopital
% source: exams/מבחנים/תשעח/מועד ב תשעח.pdf

\begin{question}
חשבו את הגבול הבא:
גבול הפונקציה $\displaystyle\lim_{x\to1^-}\frac{\frac\pi2-\arcsin x}{\sqrt[3]{1-x^2}}$.
\end{question}

\begin{hint}
זהו ביטוי מהצורה $\frac00$; הפעילו את כלל לופיטל ופשטו את מנת הנגזרות.
\end{hint}

\begin{solution}
כאשר $x\to1^-$: $\arcsin x\to\frac\pi2$ (רציפות), ולכן המונה שואף ל-$0$; גם המכנה $\sqrt[3]{1-x^2}\to0$. זהו ביטוי מהצורה $\frac00$. בקטע $(0,1)$ שתי הפונקציות גזירות:
\[ \left(\frac\pi2-\arcsin x\right)'=-\frac{1}{\sqrt{1-x^2}},\qquad \left((1-x^2)^{1/3}\right)'=\frac13(1-x^2)^{-2/3}\cdot(-2x)=-\frac{2x}{3(1-x^2)^{2/3}}\neq0. \]
מנת הנגזרות:
\[ \frac{-\frac{1}{(1-x^2)^{1/2}}}{-\frac{2x}{3(1-x^2)^{2/3}}}=\frac{3(1-x^2)^{2/3}}{2x(1-x^2)^{1/2}}=\frac{3(1-x^2)^{1/6}}{2x}\xrightarrow[x\to1^-]{}\frac{3\cdot0}{2}=0. \]
גבול מנת הנגזרות קיים, ולכן לפי כלל לופיטל
\[ \boxed{\lim_{x\to1^-}\frac{\frac\pi2-\arcsin x}{\sqrt[3]{1-x^2}}=0} \]
(אינטואיציה: $\frac\pi2-\arcsin x=\arccos x\approx\sqrt{2(1-x)}$, שהוא "גדול פחות" מ-$\sqrt[3]{1-x^2}\approx\sqrt[3]{2(1-x)}$.)
\end{solution}
